\documentclass[11pt]{article}

% Page layout
\usepackage[a4paper, margin=1in]{geometry}
\usepackage{lineno}
\linenumbers

% Typography and math
\usepackage{amsmath,amssymb}
\usepackage{microtype}

% Tables and figures
\usepackage{booktabs}
\usepackage{graphicx}
\usepackage{multirow}
\usepackage{caption}

% References and links
\usepackage[colorlinks=true,linkcolor=blue,citecolor=blue,urlcolor=blue]{hyperref}
\usepackage{xcolor}

% Code listings
\usepackage{listings}
\lstset{
    basicstyle=\small\ttfamily,
    frame=single,
    backgroundcolor=\color{gray!5},
    breaklines=true,
    columns=fullflexible,
}

% Algorithms
\usepackage{algorithm}
\usepackage{algorithmic}

\title{Supplementary Materials:\\Anchor Transfer Learning for Cross-Dataset\\Drug-Target Affinity Prediction}

\author{Ba\c{s}ar Temiz\\Bo\u{g}azi\c{c}i University, Computer Engineering Department}

\date{}

\begin{document}

\maketitle

\tableofcontents
\newpage

%% ================================================================
\section{Model Implementation Details}
\label{sec:model_details}
%% ================================================================

This section provides the complete architectural specification for each model evaluated in the paper. All models share the same training infrastructure and data splits for fair comparison. Source code is available at \url{https://github.com/Basartemiz/IDR-GAT}.

\subsection{Anchor Transfer V2 (Primary Model)}

Anchor Transfer V2 is the primary proposed architecture. It receives a triplet $(a, q, d)$ consisting of an anchor protein, a query protein, and a drug, and predicts both a binary binding label and a continuous pKi value.

\paragraph{Protein encoder.}
Both anchor and query proteins are represented by frozen ESM-2 embeddings. We evaluate two ESM-2 scales:
\begin{itemize}
    \item \textbf{V2-35M}: ESM-2 \texttt{esm2\_t12\_35M\_UR50D}, representation dimension 480, extracted at layer~12.
    \item \textbf{V2-650M}: ESM-2 \texttt{esm2\_t33\_650M\_UR50D}, representation dimension 1280, extracted at layer~33.
\end{itemize}
The ESM-2 parameters are kept frozen throughout training. A shared projection module maps both anchor and query embeddings to a 256-dimensional task-specific space:
\[
h_p = \mathrm{LayerNorm}\!\left(\mathrm{ReLU}\!\left(W_p\, x_p + b_p\right)\right), \quad h_p \in \mathbb{R}^{256}.
\]
Weight sharing ensures that anchor and query proteins are directly comparable in the same latent space.

\paragraph{Drug encoder.}
Drugs are encoded from SMILES strings using a parallel convolutional encoder. SMILES characters are mapped to integer indices using a 52-token vocabulary (Table~\ref{tab:smiles_vocab}) and embedded into 128-dimensional vectors. Three parallel 1D convolution branches with kernel sizes 4, 6, and 8 (each with 32 filters) are applied, followed by global max pooling over the sequence dimension. The pooled outputs are concatenated into a 96-dimensional feature vector and projected to 256 dimensions via a linear layer:
\[
h_d = W_d\,\big[\mathrm{pool}(\mathrm{Conv}_4) \,\|\, \mathrm{pool}(\mathrm{Conv}_6) \,\|\, \mathrm{pool}(\mathrm{Conv}_8)\big] + b_d, \quad h_d \in \mathbb{R}^{256}.
\]
SMILES strings are truncated or right-padded to a maximum length of 100 characters.

\paragraph{Triple pairwise interaction module.}
Each of the three pairwise interaction blocks takes a 512-dimensional concatenation of two entity embeddings and produces a 256-dimensional interaction feature. Each block is a two-layer MLP:
\[
\phi(x, y) = \mathrm{ReLU}\!\Big(W_2\,\mathrm{ReLU}\!\big(\mathrm{Dropout}_{0.3}(W_1\,[x \| y] + b_1)\big) + b_2\Big), \quad \phi \in \mathbb{R}^{256}.
\]
The three interaction features are:
\begin{itemize}
    \item $h_{ad} = \phi_{ad}(h_a, h_d)$: anchor--drug interaction (known binding evidence).
    \item $h_{qd} = \phi_{qd}(h_q, h_d)$: query--drug interaction (direct compatibility).
    \item $h_{aq} = \phi_{aq}(h_a, h_q)$: anchor--query interaction (protein similarity).
\end{itemize}

\paragraph{Fusion and prediction heads.}
The six features are concatenated into a 1536-dimensional fused representation:
\[
z = [h_a \| h_q \| h_d \| h_{ad} \| h_{qd} \| h_{aq}] \in \mathbb{R}^{1536}.
\]
Two task-specific heads process this representation:
\begin{itemize}
    \item \textbf{Classifier}: Linear(1536$\to$512) $\to$ ReLU $\to$ Dropout(0.3) $\to$ Linear(512$\to$256) $\to$ ReLU $\to$ Dropout(0.3) $\to$ Linear(256$\to$1) $\to$ Sigmoid.
    \item \textbf{Regressor}: Same architecture as the classifier, but with a linear output instead of sigmoid.
\end{itemize}

\paragraph{Parameter count.}
Table~\ref{tab:param_counts} reports the number of trainable parameters for each model. ESM-2 parameters are excluded as they are frozen.

% ---------------------------------------------------------------
\subsection{DeepDTA Baseline}

Our DeepDTA implementation follows \"{O}zt\"{u}rk et al.\ (2018) with minor adaptations for our training infrastructure.

\paragraph{Protein encoder.}
Protein sequences are encoded at the character level using a 26-token amino acid vocabulary ($\{$A--Z$\}$). Sequences are truncated or right-padded to a maximum length of 1000 residues. A learnable embedding layer maps each residue to a 128-dimensional vector. Three sequential 1D convolutions with 32, 64, and 96 filters (kernel size~8) are applied, followed by global max pooling, producing a 96-dimensional protein feature.

\paragraph{Drug encoder.}
SMILES strings are encoded using a 66-token vocabulary (covering digits, operators, and atom symbols) with maximum length 100. The same convolutional architecture as the protein encoder is applied: embedding dimension 128, three sequential convolutions (32$\to$64$\to$96 filters, kernel size~8), and global max pooling to produce a 96-dimensional drug feature.

\paragraph{Prediction head.}
The protein and drug features are concatenated into a 192-dimensional vector and passed through a four-layer MLP: Linear(192$\to$1024) $\to$ ReLU $\to$ Dropout(0.1) $\to$ Linear(1024$\to$1024) $\to$ ReLU $\to$ Dropout(0.1) $\to$ Linear(1024$\to$512) $\to$ ReLU $\to$ Dropout(0.1) $\to$ Linear(512$\to$1).

\paragraph{Training.}
DeepDTA is trained with MSE loss only (no binary classification head). All other training hyperparameters match the shared protocol in Section~\ref{sec:training}.

% ---------------------------------------------------------------
\subsection{ConPlex* Baseline}

ConPlex* is our re-implementation of the ConPlex contrastive co-embedding approach (Singh et al., 2023), adapted to use ESM-2 embeddings and a SMILES convolutional encoder instead of the original ProtBERT and Morgan fingerprints. This modification isolates the effect of the contrastive co-embedding formulation from the choice of protein and drug encoders.

\paragraph{Protein projector.}
Frozen ESM-2 35M embeddings (480-dim) are projected through a two-layer MLP: Linear(480$\to$512) $\to$ ReLU $\to$ Dropout(0.1) $\to$ Linear(512$\to$1024), followed by L2 normalization to unit length.

\paragraph{Drug projector.}
SMILES strings are encoded with the same parallel CNN used in the V2 model (output dimension 256), then projected: Linear(256$\to$512) $\to$ ReLU $\to$ Dropout(0.1) $\to$ Linear(512$\to$1024), followed by L2 normalization.

\paragraph{Prediction.}
Binding affinity is predicted as the cosine similarity between L2-normalized protein and drug embeddings:
\[
\mathrm{score}(p, d) = \frac{h_p \cdot h_d}{\|h_p\|\,\|h_d\|} = h_p \cdot h_d,
\]
where the second equality holds because both vectors are pre-normalized.

\paragraph{Training.}
ConPlex* alternates between two training phases each epoch:
\begin{itemize}
    \item \textbf{BCE phase} (odd epochs): The cosine similarity is rescaled to $[0, 1]$ via $\mathrm{prob} = (\mathrm{score} + 1)/2$, and binary cross-entropy is computed against the binding label.
    \item \textbf{Contrastive phase} (even epochs): A triplet margin loss is computed with margin $m = 0.5$:
    \[
    \mathcal{L}_{\mathrm{triplet}} = \mathrm{mean}\!\Big(\mathrm{ReLU}\!\big((1 - s^+) - (1 - s^-) + m\big)\Big),
    \]
    where $s^+$ is the cosine similarity for the positive drug and $s^-$ for a randomly sampled negative drug.
\end{itemize}

% ---------------------------------------------------------------
\subsection{ESM-DTA Baseline}

ESM-DTA replaces DeepDTA's learned protein character embedding with frozen ESM-2 features while keeping all other components identical. This baseline tests whether stronger pretrained protein representations alone account for the improvement, without any anchor mechanism.

\paragraph{Architecture.}
Frozen ESM-2 35M embeddings (480-dim) are projected to 256 dimensions via a linear layer with ReLU activation. The drug encoder is the same SMILES CNN used in the V2 model. Protein and drug embeddings are concatenated and passed through an MLP prediction head identical to DeepDTA's.

\paragraph{Training.}
ESM-DTA is trained with MSE loss only (no binary classification head), matching DeepDTA's supervision signal. All other hyperparameters follow the shared protocol.

% ---------------------------------------------------------------
\subsection{Drug-Anchor Ablation}

The Drug-Anchor ablation reverses the anchor modality: instead of conditioning on a known protein binder, it conditions on a known drug binder. The triplet becomes $(a_d, q_d, p)$, where $a_d$ is an anchor drug (the drug with the highest observed pKi for protein $p$), $q_d$ is the query drug, and $p$ is the target protein.

\paragraph{Architecture.}
The protein branch uses the same frozen ESM-2 projection as V2. Two instances of the SMILES CNN (with shared weights) encode the anchor drug and query drug. Three pairwise interaction blocks model $(a_d, p)$, $(q_d, p)$, and $(a_d, q_d)$ interactions, mirroring the V2 structure. Fusion and dual prediction heads are identical to V2.

\paragraph{Anchor selection.}
For each protein in the training set, the anchor drug is the compound with the highest observed pKi. When the query drug is itself the top-ranked compound, the second-highest pKi drug is used as the anchor.

%% ================================================================
\section{Training Procedure}
\label{sec:training}
%% ================================================================

\subsection{Shared Hyperparameters}

All models were trained with the following shared hyperparameters unless otherwise noted:

\begin{table}[h]
\centering
\caption{Shared training hyperparameters across all models.}
\label{tab:hyperparams}
\begin{tabular}{ll}
\toprule
Hyperparameter & Value \\
\midrule
Optimizer & AdamW \\
Learning rate & $10^{-3}$ \\
Weight decay & $10^{-4}$ \\
LR schedule & Cosine annealing ($T_{\max} = \mathrm{epochs}$) \\
Batch size & 256 \\
Maximum epochs & 100 \\
Early stopping patience & 20 (on validation loss) \\
Gradient clipping & Max norm 1.0 \\
Random seed & 42 \\
\bottomrule
\end{tabular}
\end{table}

\subsection{Multi-Task Loss}

Models with dual prediction heads (V2-35M, V2-650M, Drug-Anchor) are trained with a combined loss:
\[
\mathcal{L} = \mathcal{L}_{\mathrm{BCE}} + \alpha\,\mathcal{L}_{\mathrm{MSE}}, \quad \alpha = 1.0.
\]
The MSE term is computed over all training examples. The BCE term is computed only over samples with confident labels (pKi $\geq 7$ or pKi $\leq 5$); samples in the ambiguous range $5 < \mathrm{pKi} < 7$ are excluded from binary supervision but retained for regression. This masking avoids noisy gradient signal near the decision boundary.

Pairwise baselines (DeepDTA, ESM-DTA) are trained with MSE loss only. ConPlex* uses alternating BCE and triplet losses as described above.

\subsection{Data Splitting Protocol}

All datasets are split at the protein level to prevent information leakage:
\begin{enumerate}
    \item Collect all unique protein identifiers (UniProt IDs).
    \item Sort and shuffle with seed 42.
    \item Assign the first 10\% to the test set, the next 10\% to the validation set, and the remaining 80\% to training.
    \item All interactions involving a given protein are assigned to the same split.
\end{enumerate}
This ensures that no protein appears in both training and evaluation, which is stricter than random interaction-level splitting.

\subsection{Checkpoint Selection}

The model checkpoint with the lowest validation loss is saved during training. Early stopping terminates training if validation loss does not improve for 20 consecutive epochs. The best checkpoint is used for all subsequent evaluation.

%% ================================================================
\section{ESM-2 Embedding Extraction}
\label{sec:esm2}
%% ================================================================

Protein embeddings are precomputed and cached before training. We use the official ESM-2 implementation from the \texttt{fair-esm} library.

\paragraph{Model variants.}
We use two ESM-2 scales: \texttt{esm2\_t12\_35M\_UR50D} (12 layers, 480-dim representations) and \texttt{esm2\_t33\_650M\_UR50D} (33 layers, 1280-dim representations).

\paragraph{Preprocessing.}
Protein sequences are derived from UniProt entries. PDB three-letter residue codes are converted to one-letter amino acid codes. Sequences longer than 1022 residues (the ESM-2 context window, excluding BOS/EOS tokens) are truncated.

\paragraph{Extraction.}
Each sequence is tokenized with the ESM-2 alphabet, passed through the frozen model, and the representation at the final transformer layer is extracted. The output shape is $(L{+}2, D)$, where $L$ is the sequence length and $D$ is the representation dimension. BOS and EOS token representations are trimmed, and the remaining $L$ residue representations are mean-pooled to produce a single $D$-dimensional embedding per protein.

\paragraph{Batching.}
Sequences are processed in batches of 4. Sequences that exceed available GPU memory are processed individually as a fallback.

%% ================================================================
\section{Anchor Selection Protocol}
\label{sec:anchor_selection}
%% ================================================================

\subsection{Protein-Side Anchors (V2 Models)}

For each drug $d$ in the training set, the anchor protein is defined as:
\[
a(d) = \arg\max_{p \in \mathcal{P}_{\mathrm{train}}} \mathrm{pKi}(p, d),
\]
where $\mathcal{P}_{\mathrm{train}}$ is the set of training proteins with observed interactions for drug $d$. When the query protein $q$ is itself the top-ranked binder, the anchor falls back to the protein with the second-highest pKi to ensure $a \neq q$.

At inference on external benchmarks, anchors are selected in oracle mode: for each query interaction $(q, d)$, the anchor is chosen from known strong binders within the benchmark itself. This tests the model's ability to transfer from an observed binding context, decoupled from anchor retrieval quality.

\subsection{Drug-Side Anchors (Drug-Anchor Ablation)}

For each protein $p$, the anchor drug is:
\[
a_d(p) = \arg\max_{d \in \mathcal{D}_{\mathrm{train}}} \mathrm{pKi}(p, d),
\]
with the same fallback logic when the query drug equals the top-ranked compound.

%% ================================================================
\section{Dataset Construction Details}
\label{sec:datasets}
%% ================================================================

\subsection{Drug Target Commons (DTC)}

The primary training dataset was obtained from Drug Target Commons. We retained only interactions with Ki-derived pKi values to maintain a consistent inhibition-constant scale. pKi values are computed as:
\[
\mathrm{pKi} = 9 - \log_{10}(\mathrm{Ki}_{\mathrm{nM}}),
\]
where $\mathrm{Ki}_{\mathrm{nM}}$ is the inhibition constant in nanomolar units. Values outside the range $[3.0, 12.0]$ are discarded as likely measurement artifacts. The final dataset contains 401,978 interactions spanning 2,827 proteins and 164,831 drugs.

\subsection{Davis Kinase Benchmark}

The Davis dataset reports dissociation constants (Kd) rather than inhibition constants (Ki). We treat pKd and pKi interchangeably for ranking evaluation, as both measure binding strength on comparable logarithmic scales. This is a common simplification in the DTA literature. The dataset contains 30,056 kinase--drug interactions over 442 proteins and 68 drugs. Approximately 70\% of values are censored at pKi~$= 5$ (corresponding to Kd $> 10{,}000$~nM), making the binary classification task highly imbalanced.

\subsection{GLASS GPCR Benchmark}

The GLASS 2.0 dataset was filtered to retain only Ki measurements. We removed all proteins overlapping with DTC, yielding 16,721 interactions across 67 novel GPCR proteins. This benchmark probes family-level transfer: the training data is kinase-dominated, whereas GPCRs have different sequence patterns and binding landscapes.

\subsection{BDB-Ki Clean Benchmark}

The BDB-Ki clean benchmark was curated from BindingDB by removing all proteins and drugs that overlap with DTC. The final dataset contains 37,588 interactions involving 324 proteins and 24,662 drugs. This is the strictest test: neither side of the interaction pair is shared with the source dataset.

\subsection{IDP Benchmark}

The IDP benchmark contains 176,840 interactions across 203 proteins, including 170 intrinsically disordered proteins and 33 ordered proteins included for comparison. Relative to DTC, the benchmark has 0\% protein overlap, 0\% interaction overlap, and 1.9\% residual drug overlap. IDPs were identified using DisProt annotations and literature curation. This benchmark targets the most challenging transfer regime, as IDPs lack the stable sequence--structure regularities that standard affinity predictors rely on.

\subsection{SMILES Encoding Vocabulary}

\begin{table}[h]
\centering
\caption{SMILES character vocabulary used across all models (52 tokens). Index 0 is reserved for padding.}
\label{tab:smiles_vocab}
\small
\begin{tabular}{ll}
\toprule
Category & Characters \\
\midrule
Structural & \texttt{( ) [ ] . = \# \% + - / @} \\
Digits & \texttt{0 1 2 3 4 5 6 7 8 9} \\
Uppercase atoms & \texttt{A B C D E F G H I K L M N O P R S T V Y} \\
Lowercase atoms & \texttt{c e l n o r s t u} \\
\bottomrule
\end{tabular}
\end{table}

All SMILES strings are truncated or right-padded to 100 characters. Characters not present in the vocabulary are mapped to index~0 (padding). DeepDTA uses a separate 66-token vocabulary that follows the original implementation.

%% ================================================================
\section{Evaluation Protocol Details}
\label{sec:eval_details}
%% ================================================================

\subsection{Common Anchored Subsets}

To ensure fair comparison across methods, all models are evaluated on the same subset of interactions for each benchmark. This \emph{common anchored subset} is defined as the set of interactions for which a valid anchor can be constructed for all anchor-based methods. Pairwise baselines (DeepDTA, ESM-DTA, ConPlex*) are evaluated on the same subset even though they do not use anchors, so that sample sizes are identical across methods.

\subsection{Metrics}

We report three complementary metrics:

\paragraph{AUROC.}
Area Under the Receiver Operating Characteristic curve, computed on binary labels derived from pKi thresholds: positive if pKi $\geq 7$, negative if pKi $\leq 5$. Interactions in the ambiguous range ($5 < \mathrm{pKi} < 7$) are excluded from AUROC computation.

\paragraph{Concordance Index (CI).}
The fraction of correctly ordered pairs among all pairs with different observed pKi values:
\[
\mathrm{CI} = \frac{\sum_{i<j} \mathbf{1}[\hat{y}_i > \hat{y}_j \text{ and } y_i > y_j] + 0.5 \cdot \mathbf{1}[\hat{y}_i = \hat{y}_j \text{ and } y_i > y_j]}{\sum_{i<j} \mathbf{1}[y_i \neq y_j]}.
\]
When the number of pairs exceeds 100,000, we subsample uniformly at random for computational efficiency.

\paragraph{RMSE.}
Root mean squared error between predicted and observed pKi values, computed over all interactions (including the ambiguous range).

\subsection{Oracle Anchor Mode}

External benchmarks are evaluated in oracle anchor mode. For each query interaction $(q, d)$, the anchor is selected as the protein (or drug, for Drug-Anchor) with the highest observed pKi for the same drug (or protein) \emph{within the benchmark dataset itself}. This decouples the evaluation of the transfer mechanism from the quality of anchor retrieval. In a prospective deployment, anchor retrieval would introduce additional error; our oracle evaluation therefore represents an upper bound on performance.

%% ================================================================
\section{Conformation Weighting Analysis}
\label{sec:conformation}
%% ================================================================

As an auxiliary negative-control experiment, we tested whether replacing the query protein's ESM-2 embedding with a structurally weighted average of ordered-protein embeddings would improve transfer to IDPs. The rationale was that disordered proteins might benefit from borrowing representation context from structurally similar ordered proteins.

\paragraph{Procedure.}
For each IDP query, we used Foldseek to search the IDRome database of predicted IDP conformations and retrieved structurally matched ordered proteins. The query embedding was replaced with a TM-score-weighted average of the retrieved ordered-protein ESM-2 embeddings. The rest of the model and evaluation protocol remained unchanged.

\paragraph{Result.}
On the common evaluation subset ($n = 2{,}442$), the standard V2 model reached AUROC 0.584, whereas the conformation-weighted variant reached 0.575. Structural similarity therefore does not appear to carry the binding information needed for transfer in this setting. The anchor mechanism helps because it transfers functional knowledge (observed binding evidence), not structural proximity.

%% ================================================================
\section{Model Parameter Counts}
\label{sec:params}
%% ================================================================

\begin{table}[h]
\centering
\caption{Trainable parameter counts for each model. ESM-2 parameters are frozen and excluded from the count.}
\label{tab:param_counts}
\begin{tabular}{lr}
\toprule
Model & Trainable Parameters \\
\midrule
V2-35M & 2.36M \\
V2-650M & 2.62M \\
Drug-Anchor & 2.36M \\
DeepDTA & 1.19M \\
ESM-DTA & 1.05M \\
ConPlex* & 1.60M \\
\bottomrule
\end{tabular}
\end{table}

The anchor transfer models have approximately twice the parameters of pairwise baselines due to the triple pairwise interaction module and dual prediction heads. Despite this, the pairwise baselines achieve higher same-dataset performance (Table~2 in the main paper), indicating that the anchor transfer advantage under distribution shift is not simply a function of model capacity.

%% ================================================================
\section{Software and Hardware}
\label{sec:software}
%% ================================================================

All experiments were conducted using:
\begin{itemize}
    \item Python 3.10
    \item PyTorch 2.1
    \item ESM (\texttt{fair-esm}) for protein language model inference
    \item RDKit for SMILES parsing and molecular graph construction
    \item scikit-learn for evaluation metrics
\end{itemize}

Training was performed on NVIDIA A40 (48 GB) and A100 (80 GB) GPUs. Each model trains to convergence in approximately 30--60 minutes on DTC with a single GPU. ESM-2 embedding extraction for all DTC proteins takes approximately 15 minutes (35M) or 45 minutes (650M) on an A100.

\end{document}
